%% =====================================================================
%%  B4-T-18-01 -- what B4 contributes to "The Symbols of Authority"
%%  -------------------------------------------------------------------
%%  LAYER A: APPEND-ONLY. Written once, then left alone. Material found
%%  only here sits in the `unique` slot and is protected by rule R10.
%% =====================================================================
%% @kind: source
%% @id: B4-T-18-01
%% @book: B4
%% @topic: T-18-01
%% @locator: ch. 6, pp. 178--204
%% @status: distilled
%% @unique: yes
%% @integrated: yes
%% @updated: 2026-09-19

\dossier{B4}{T-18-01}{\bookshort{B4}}{ch. 6, pp. 178--204}

\begin{distilled}
  \item Compliance rises sharply with the apparent authority of the requester, independently of the merits of the request.
  \item Authority is identified by visible symbols, because judging the underlying expertise requires knowing the field.
  \item The symbols are substantially cheaper to acquire than the substance, and the response fires on the symbol without evaluation.
  \item Deference once given is difficult to withdraw, since withdrawal requires admitting the original error.
\end{distilled}

\begin{unique}
  \item The proxy account: authority symbols are used because direct assessment is unavailable, which explains both why the response is rational and why it is exploitable.
  \item The separation of the two questions --- is this a genuine expert, and is this expert being straight with me --- as the structure of the defence.
\end{unique}
